\documentclass[11pt]{article}
\usepackage[utf8]{inputenc}
\usepackage{amsmath, amssymb, amsthm}
\usepackage{mathpartir}
\usepackage{stmaryrd}
\usepackage{xcolor}
\usepackage{geometry}
\usepackage{hyperref}

\geometry{margin=1in}

\title{Vyper Typing Rules}
\author{Extracted from the Vyper Compiler Implementation}
\date{}

\newcommand{\kw}[1]{\texttt{#1}}
\newcommand{\typerule}[1]{\textsc{#1}}
\newcommand{\judge}[3]{#1 \vdash #2 : #3}
\newcommand{\judgem}[4]{#1 \vdash_{#2} #3 : #4}
\newcommand{\stmt}[4]{#1; #2 \vdash #3 \dashv #4}
\newcommand{\subtype}[2]{#1 <: #2}

\begin{document}

\maketitle

\begin{abstract}
This document presents the formal typing rules for the Vyper smart contract programming language,
extracted from the compiler implementation. Vyper is a statically typed, Pythonic language
targeting the Ethereum Virtual Machine (EVM). These rules follow standard academic notation
for type systems.
\end{abstract}

\tableofcontents
\newpage

%==============================================================================
\section{Type Syntax}
%==============================================================================

\subsection{Base Types}

\begin{align*}
\tau ::= \quad & \kw{bool} & \text{Boolean} \\
\mid \quad & \kw{int}N \quad (N \in \{8,16,\ldots,256\}) & \text{Signed integer} \\
\mid \quad & \kw{uint}N \quad (N \in \{8,16,\ldots,256\}) & \text{Unsigned integer} \\
\mid \quad & \kw{decimal} & \text{Fixed-point decimal (168 bits)} \\
\mid \quad & \kw{address} & \text{Ethereum address (160 bits)} \\
\mid \quad & \kw{bytes}M \quad (M \in \{1,\ldots,32\}) & \text{Fixed-size byte array} \\
\mid \quad & \kw{Bytes}[n] & \text{Dynamic byte array (max length $n$)} \\
\mid \quad & \kw{String}[n] & \text{Dynamic string (max length $n$)}
\end{align*}

\subsection{Composite Types}

\begin{align*}
\tau ::= \quad & \tau[n] & \text{Static array of length } n \\
\mid \quad & \kw{DynArray}[\tau, n] & \text{Dynamic array (max length } n \text{)} \\
\mid \quad & \kw{HashMap}[\kappa, \tau] & \text{Mapping from } \kappa \text{ to } \tau \\
\mid \quad & (\tau_1, \ldots, \tau_n) & \text{Tuple} \\
\mid \quad & \kw{struct}\ S & \text{User-defined struct} \\
\mid \quad & \kw{flag}\ F & \text{User-defined flag (enum)} \\
\mid \quad & \kw{event}\ E & \text{User-defined event}
\end{align*}

%==============================================================================
\section{Typing Contexts}
%==============================================================================

\subsection{Type Environment}

\[
\Gamma ::= \emptyset \mid \Gamma, x : \tau
\]

\subsection{Data Locations}

\[
\ell ::= \kw{storage} \mid \kw{transient} \mid \kw{memory} \mid \kw{calldata} \mid \kw{code}
\]

\subsection{Modifiability (for expressions and variables)}

\[
\mu ::= \kw{constant} \mid \kw{runtime\_constant} \mid \kw{modifiable}
\]

\subsection{State Mutability (for functions)}

\[
m ::= \kw{pure} \mid \kw{view} \mid \kw{nonpayable} \mid \kw{payable}
\]

With ordering: $\kw{pure} < \kw{view} < \kw{nonpayable} < \kw{payable}$

%==============================================================================
\section{Subtyping Rules}
%==============================================================================

\fbox{$\subtype{\tau_1}{\tau_2}$} \quad ($\tau_1$ is a subtype of $\tau_2$)

\subsection{Structural Rules}

\begin{mathpar}
\inferrule*[right=S-Refl]
  { }
  {\subtype{\tau}{\tau}}

\inferrule*[right=S-Trans]
  {\subtype{\tau_1}{\tau_2} \\ \subtype{\tau_2}{\tau_3}}
  {\subtype{\tau_1}{\tau_3}}
\end{mathpar}

\subsection{Numeric Types (No Implicit Coercion)}

\begin{mathpar}
\inferrule*[right=S-Int]
  {N_1 = N_2 \\ s_1 = s_2}
  {\subtype{\kw{int}^{s_1}_{N_1}}{\kw{int}^{s_2}_{N_2}}}
\end{mathpar}

\subsection{Array Subtyping}

\begin{mathpar}
\inferrule*[right=S-SArray]
  {\subtype{\tau_1}{\tau_2} \\ \subtype{\tau_2}{\tau_1} \\ n_1 = n_2}
  {\subtype{\tau_1[n_1]}{\tau_2[n_2]}}

\inferrule*[right=S-DynArray]
  {\subtype{\tau_1}{\tau_2} \\ n_1 \leq n_2}
  {\subtype{\kw{DynArray}[\tau_1, n_1]}{\kw{DynArray}[\tau_2, n_2]}}
\end{mathpar}

\subsection{Bytestring Subtyping}

\begin{mathpar}
\inferrule*[right=S-Bytes]
  {n_1 \leq n_2}
  {\subtype{\kw{Bytes}[n_1]}{\kw{Bytes}[n_2]}}

\inferrule*[right=S-String]
  {n_1 \leq n_2}
  {\subtype{\kw{String}[n_1]}{\kw{String}[n_2]}}
\end{mathpar}

\subsection{Special Types}

\begin{mathpar}
\inferrule*[right=S-Self]
  { }
  {\subtype{\kw{self}}{\kw{address}}}

\inferrule*[right=S-Struct]
  {S_1 \equiv S_2}
  {\subtype{\kw{struct}\ S_1}{\kw{struct}\ S_2}}

\inferrule*[right=S-Flag]
  {F_1 \equiv F_2}
  {\subtype{\kw{flag}\ F_1}{\kw{flag}\ F_2}}
\end{mathpar}

%==============================================================================
\section{Expression Typing Rules}
%==============================================================================

\fbox{$\judge{\Gamma}{e}{\tau}$} \quad (expression $e$ has type $\tau$ in context $\Gamma$)

\subsection{Literals}

\begin{mathpar}
\inferrule*[right=T-BoolTrue]
  { }
  {\judge{\Gamma}{\kw{True}}{\kw{bool}}}

\inferrule*[right=T-BoolFalse]
  { }
  {\judge{\Gamma}{\kw{False}}{\kw{bool}}}

\inferrule*[right=T-IntLit]
  {-2^{N-1} \leq n \leq 2^{N-1} - 1}
  {\judge{\Gamma}{n}{\kw{int}N}}

\inferrule*[right=T-UintLit]
  {0 \leq n \leq 2^N - 1}
  {\judge{\Gamma}{n}{\kw{uint}N}}

\inferrule*[right=T-DecLit]
  {\text{lo} \leq d \leq \text{hi}}
  {\judge{\Gamma}{d}{\kw{decimal}}}

\inferrule*[right=T-Addr]
  {\text{addr is valid checksummed address}}
  {\judge{\Gamma}{\kw{0x}\mathit{addr}}{\kw{address}}}

\inferrule*[right=T-BytesM]
  {|s| = M \\ 1 \leq M \leq 32}
  {\judge{\Gamma}{\kw{0x}s}{\kw{bytes}M}}

\inferrule*[right=T-BytesLit]
  {|s| \leq n}
  {\judge{\Gamma}{\kw{b"}s\kw{"}}{\kw{Bytes}[n]}}

\inferrule*[right=T-StrLit]
  {|s| \leq n}
  {\judge{\Gamma}{\kw{"}s\kw{"}}{\kw{String}[n]}}
\end{mathpar}

\subsection{Variables}

\begin{mathpar}
\inferrule*[right=T-Var]
  {x : \tau \in \Gamma}
  {\judge{\Gamma}{x}{\tau}}
\end{mathpar}

\subsection{Boolean Operations}

\begin{mathpar}
\inferrule*[right=T-And]
  {\judge{\Gamma}{e_1}{\kw{bool}} \\ \judge{\Gamma}{e_2}{\kw{bool}}}
  {\judge{\Gamma}{e_1\ \kw{and}\ e_2}{\kw{bool}}}

\inferrule*[right=T-Or]
  {\judge{\Gamma}{e_1}{\kw{bool}} \\ \judge{\Gamma}{e_2}{\kw{bool}}}
  {\judge{\Gamma}{e_1\ \kw{or}\ e_2}{\kw{bool}}}

\inferrule*[right=T-Not]
  {\judge{\Gamma}{e}{\kw{bool}}}
  {\judge{\Gamma}{\kw{not}\ e}{\kw{bool}}}
\end{mathpar}

\subsection{Arithmetic Operations}

Let $\mathit{NumericT} = \{\kw{int}N, \kw{uint}N, \kw{decimal}\}$

\begin{mathpar}
\inferrule*[right=T-Arith]
  {\judge{\Gamma}{e_1}{\tau} \\ \judge{\Gamma}{e_2}{\tau} \\ \tau \in \mathit{NumericT} \\
   \oplus \in \{+, -, *, /\!/, \%\}}
  {\judge{\Gamma}{e_1 \oplus e_2}{\tau}}

\inferrule*[right=T-DecDiv]
  {\judge{\Gamma}{e_1}{\kw{decimal}} \\ \judge{\Gamma}{e_2}{\kw{decimal}}}
  {\judge{\Gamma}{e_1\ /\ e_2}{\kw{decimal}}}

\inferrule*[right=T-Neg]
  {\judge{\Gamma}{e}{\tau} \\ \tau \in \{\kw{int}N, \kw{decimal}\}}
  {\judge{\Gamma}{-e}{\tau}}
\end{mathpar}

\subsection{Exponentiation}

\begin{mathpar}
\inferrule*[right=T-Pow]
  {\judge{\Gamma}{e_1}{\tau} \\ \judge{\Gamma}{e_2}{\tau} \\
   \tau \in \mathit{IntegerT} \\ \text{one of } e_1, e_2 \text{ is literal}}
  {\judge{\Gamma}{e_1\ \kw{**}\ e_2}{\tau}}
\end{mathpar}

\subsection{Bitwise Operations}

Let $\mathit{BitwiseT} = \{\kw{uint}N, \kw{bytes}M, \kw{flag}\}$

\begin{mathpar}
\inferrule*[right=T-Bitwise]
  {\judge{\Gamma}{e_1}{\tau} \\ \judge{\Gamma}{e_2}{\tau} \\
   \tau \in \mathit{BitwiseT} \\ \odot \in \{\&, |, \wedge\}}
  {\judge{\Gamma}{e_1\ \odot\ e_2}{\tau}}

\inferrule*[right=T-BitNot]
  {\judge{\Gamma}{e}{\tau} \\ \tau \in \{\kw{uint256}, \kw{bytes32}, \kw{flag}\}}
  {\judge{\Gamma}{\sim e}{\tau}}
\end{mathpar}

\subsection{Shift Operations (256-bit types only)}

\begin{mathpar}
\inferrule*[right=T-LShift]
  {\judge{\Gamma}{e_1}{\tau} \\ \judge{\Gamma}{e_2}{\kw{uint}N} \\
   \tau \in \{\kw{uint256}, \kw{int256}, \kw{bytes32}\}}
  {\judge{\Gamma}{e_1 \ll e_2}{\tau}}

\inferrule*[right=T-RShift]
  {\judge{\Gamma}{e_1}{\tau} \\ \judge{\Gamma}{e_2}{\kw{uint}N} \\
   \tau \in \{\kw{uint256}, \kw{int256}, \kw{bytes32}\}}
  {\judge{\Gamma}{e_1 \gg e_2}{\tau}}
\end{mathpar}

\subsection{Comparisons}

\begin{mathpar}
\inferrule*[right=T-NumCmp]
  {\judge{\Gamma}{e_1}{\tau} \\ \judge{\Gamma}{e_2}{\tau} \\
   \tau \in \mathit{NumericT} \\ \bowtie\ \in \{<, \leq, >, \geq, ==, !=\}}
  {\judge{\Gamma}{e_1\ \bowtie\ e_2}{\kw{bool}}}

\inferrule*[right=T-Eq]
  {\judge{\Gamma}{e_1}{\tau} \\ \judge{\Gamma}{e_2}{\tau}}
  {\judge{\Gamma}{e_1 == e_2}{\kw{bool}}}

\inferrule*[right=T-NEq]
  {\judge{\Gamma}{e_1}{\tau} \\ \judge{\Gamma}{e_2}{\tau}}
  {\judge{\Gamma}{e_1\ !=\ e_2}{\kw{bool}}}
\end{mathpar}

\subsection{Membership Tests}

\begin{mathpar}
\inferrule*[right=T-InSArray]
  {\judge{\Gamma}{e}{\tau} \\ \judge{\Gamma}{a}{\tau[n]}}
  {\judge{\Gamma}{e\ \kw{in}\ a}{\kw{bool}}}

\inferrule*[right=T-InDArray]
  {\judge{\Gamma}{e}{\tau} \\ \judge{\Gamma}{a}{\kw{DynArray}[\tau, n]}}
  {\judge{\Gamma}{e\ \kw{in}\ a}{\kw{bool}}}

\inferrule*[right=T-InFlag]
  {\judge{\Gamma}{e}{\kw{flag}\ F} \\ \judge{\Gamma}{f}{\kw{flag}\ F}}
  {\judge{\Gamma}{e\ \kw{in}\ f}{\kw{bool}}}
\end{mathpar}

\subsection{Subscript/Index Access}

\begin{mathpar}
\inferrule*[right=T-SArrayIdx]
  {\judge{\Gamma}{e}{\tau[n]} \\ \judge{\Gamma}{i}{\mathit{IntegerT}} \\ 0 \leq i < n}
  {\judge{\Gamma}{e[i]}{\tau}}

\inferrule*[right=T-DArrayIdx]
  {\judge{\Gamma}{e}{\kw{DynArray}[\tau, n]} \\ \judge{\Gamma}{i}{\mathit{IntegerT}}}
  {\judge{\Gamma}{e[i]}{\tau}}

\inferrule*[right=T-MapIdx]
  {\judge{\Gamma}{m}{\kw{HashMap}[\kappa, \tau]} \\ \judge{\Gamma}{k}{\kappa}}
  {\judge{\Gamma}{m[k]}{\tau}}

\inferrule*[right=T-TupleIdx]
  {\judge{\Gamma}{t}{(\tau_0, \ldots, \tau_{n-1})} \\ 0 \leq i < n}
  {\judge{\Gamma}{t[i]}{\tau_i}}
\end{mathpar}

\subsection{Attribute Access}

\begin{mathpar}
\inferrule*[right=T-StructAttr]
  {\judge{\Gamma}{e}{\kw{struct}\ S} \\ f : \tau \in S}
  {\judge{\Gamma}{e.f}{\tau}}

\inferrule*[right=T-AddrBal]
  {\judge{\Gamma}{e}{\kw{address}}}
  {\judge{\Gamma}{e.\kw{balance}}{\kw{uint256}}}

\inferrule*[right=T-AddrSize]
  {\judge{\Gamma}{e}{\kw{address}}}
  {\judge{\Gamma}{e.\kw{codesize}}{\kw{uint256}}}

\inferrule*[right=T-AddrHash]
  {\judge{\Gamma}{e}{\kw{address}}}
  {\judge{\Gamma}{e.\kw{codehash}}{\kw{bytes32}}}

\inferrule*[right=T-AddrContract]
  {\judge{\Gamma}{e}{\kw{address}}}
  {\judge{\Gamma}{e.\kw{is\_contract}}{\kw{bool}}}
\end{mathpar}

\subsection{Composite Literals}

\begin{mathpar}
\inferrule*[right=T-SArrayLit]
  {\judge{\Gamma}{e_1}{\tau} \\ \cdots \\ \judge{\Gamma}{e_n}{\tau}}
  {\judge{\Gamma}{[e_1, \ldots, e_n]}{\tau[n]}}

\inferrule*[right=T-DArrayLit]
  {\judge{\Gamma}{e_1}{\tau} \\ \cdots \\ \judge{\Gamma}{e_n}{\tau}}
  {\judge{\Gamma}{[e_1, \ldots, e_n]}{\kw{DynArray}[\tau, n]}}

\inferrule*[right=T-Tuple]
  {\judge{\Gamma}{e_1}{\tau_1} \\ \cdots \\ \judge{\Gamma}{e_n}{\tau_n}}
  {\judge{\Gamma}{(e_1, \ldots, e_n)}{(\tau_1, \ldots, \tau_n)}}
\end{mathpar}

\subsection{Struct Construction}

\begin{mathpar}
\inferrule*[right=T-StructCtor]
  {\kw{struct}\ S = \{f_1 : \tau_1, \ldots, f_n : \tau_n\} \\
   \judge{\Gamma}{e_i}{\tau_i} \text{ for all } i}
  {\judge{\Gamma}{S(f_1 = e_1, \ldots, f_n = e_n)}{\kw{struct}\ S}}
\end{mathpar}

\subsection{Conditional Expression}

\begin{mathpar}
\inferrule*[right=T-IfExp]
  {\judge{\Gamma}{e_c}{\kw{bool}} \\ \judge{\Gamma}{e_1}{\tau} \\ \judge{\Gamma}{e_2}{\tau}}
  {\judge{\Gamma}{e_1\ \kw{if}\ e_c\ \kw{else}\ e_2}{\tau}}
\end{mathpar}

\subsection{Type Conversion}

\begin{mathpar}
\inferrule*[right=T-Convert]
  {\judge{\Gamma}{e}{\tau_1} \\ \text{convert}(\tau_1, \tau_2) \text{ is valid}}
  {\judge{\Gamma}{\kw{convert}(e, \tau_2)}{\tau_2}}
\end{mathpar}

\subsection{Dynamic Array Methods}

\begin{mathpar}
\inferrule*[right=T-Append]
  {\judge{\Gamma}{a}{\kw{DynArray}[\tau, n]} \\ \judge{\Gamma}{e}{\tau}}
  {\judge{\Gamma}{a.\kw{append}(e)}{()}}

\inferrule*[right=T-Pop]
  {\judge{\Gamma}{a}{\kw{DynArray}[\tau, n]}}
  {\judge{\Gamma}{a.\kw{pop}()}{\tau}}
\end{mathpar}

%==============================================================================
\section{Function Typing Rules}
%==============================================================================

\fbox{$\judgem{\Gamma}{m}{f}{(\overline{\tau}) \to \tau_r}$} \quad (function $f$ with mutability $m$)

\subsection{Function Definition}

\begin{mathpar}
\inferrule*[right=T-FunDef]
  {\Gamma, x_1 : \tau_1, \ldots, x_n : \tau_n \vdash_m \mathit{body} \dashv \tau_r \\
   \mathit{vis} \in \{\kw{internal}, \kw{external}, \kw{deploy}\}}
  {\Gamma \vdash \kw{@}\mathit{vis}\ \kw{@}m\ \kw{def}\ f(x_1 : \tau_1, \ldots, x_n : \tau_n) \to \tau_r :
   (\tau_1, \ldots, \tau_n) \to \tau_r}
\end{mathpar}

\subsection{Function Calls}

\begin{mathpar}
\inferrule*[right=T-InternalCall]
  {\judge{\Gamma}{f}{(\tau_1, \ldots, \tau_n) \to \tau_r} \\
   \judge{\Gamma}{e_i}{\tau_i'} \\ \subtype{\tau_i'}{\tau_i} \text{ for all } i}
  {\judge{\Gamma}{f(e_1, \ldots, e_n)}{\tau_r}}

\inferrule*[right=T-ExtCall]
  {\judge{\Gamma}{e}{\kw{interface}\ I} \\ f : (\tau_1, \ldots, \tau_n) \to \tau_r \in I \\
   \judge{\Gamma}{e_i}{\tau_i'} \\ \subtype{\tau_i'}{\tau_i} \text{ for all } i \\
   m_f \geq \kw{nonpayable}}
  {\judge{\Gamma}{\kw{extcall}\ e.f(e_1, \ldots, e_n)}{\tau_r}}

\inferrule*[right=T-StaticCall]
  {\judge{\Gamma}{e}{\kw{interface}\ I} \\ f : (\tau_1, \ldots, \tau_n) \to \tau_r \in I \\
   \judge{\Gamma}{e_i}{\tau_i'} \\ \subtype{\tau_i'}{\tau_i} \text{ for all } i \\
   m_f \in \{\kw{view}, \kw{pure}\}}
  {\judge{\Gamma}{\kw{staticcall}\ e.f(e_1, \ldots, e_n)}{\tau_r}}
\end{mathpar}

%==============================================================================
\section{Statement Typing Rules}
%==============================================================================

\fbox{$\stmt{\Gamma}{m}{s}{\Gamma'}$} \quad (statement $s$ transforms context under mutability $m$)

\subsection{Variable Declaration}

\begin{mathpar}
\inferrule*[right=T-VarDecl]
  {\judge{\Gamma}{e}{\tau'} \\ \subtype{\tau'}{\tau}}
  {\stmt{\Gamma}{m}{x : \tau = e}{\Gamma, x : \tau}}
\end{mathpar}

\subsection{Assignment}

\begin{mathpar}
\inferrule*[right=T-Assign]
  {x : \tau \in \Gamma \\ \judge{\Gamma}{e}{\tau'} \\
   \subtype{\tau'}{\tau} \\ \mu(x) = \kw{modifiable}}
  {\stmt{\Gamma}{m}{x = e}{\Gamma}}

\inferrule*[right=T-AugAssign]
  {x : \tau \in \Gamma \\ \judge{\Gamma}{e}{\tau} \\
   \tau \in \mathit{NumericT} \\ \mu(x) = \kw{modifiable}}
  {\stmt{\Gamma}{m}{x\ \oplus\kw{=}\ e}{\Gamma}}
\end{mathpar}

\subsection{Control Flow}

\begin{mathpar}
\inferrule*[right=T-If]
  {\judge{\Gamma}{e}{\kw{bool}} \\
   \stmt{\Gamma}{m}{s_1}{\Gamma_1} \\
   \stmt{\Gamma}{m}{s_2}{\Gamma_2}}
  {\stmt{\Gamma}{m}{\kw{if}\ e\kw{:}\ s_1\ \kw{else:}\ s_2}{\Gamma}}

\inferrule*[right=T-ForRange]
  {\judge{\Gamma}{e_1}{\mathit{IntegerT}} \\ \judge{\Gamma}{e_2}{\mathit{IntegerT}} \\
   \stmt{\Gamma, i : \mathit{IntegerT}}{m}{s}{\Gamma'}}
  {\stmt{\Gamma}{m}{\kw{for}\ i : \mathit{IntegerT}\ \kw{in range}(e_1, e_2)\kw{:}\ s}{\Gamma}}

\inferrule*[right=T-ForArray]
  {\judge{\Gamma}{a}{\tau[n]} \\ \stmt{\Gamma, x : \tau}{m}{s}{\Gamma'}}
  {\stmt{\Gamma}{m}{\kw{for}\ x : \tau\ \kw{in}\ a\kw{:}\ s}{\Gamma}}
\end{mathpar}

\subsection{Return, Assert, Raise}

\begin{mathpar}
\inferrule*[right=T-Return]
  {\judge{\Gamma}{e}{\tau'} \\ \subtype{\tau'}{\tau_r}}
  {\stmt{\Gamma}{m}{\kw{return}\ e}{\bot}}

\inferrule*[right=T-Assert]
  {\judge{\Gamma}{e}{\kw{bool}}}
  {\stmt{\Gamma}{m}{\kw{assert}\ e}{\Gamma}}

\inferrule*[right=T-Raise]
  {\judge{\Gamma}{e}{\kw{String}[n]} \\ n \leq 1024}
  {\stmt{\Gamma}{m}{\kw{raise}\ e}{\bot}}
\end{mathpar}

\subsection{Event Logging}

\begin{mathpar}
\inferrule*[right=T-Log]
  {\kw{event}\ E = \{a_1 : \tau_1, \ldots, a_n : \tau_n\} \\
   \judge{\Gamma}{e_i}{\tau_i} \text{ for all } i \\ m \geq \kw{nonpayable}}
  {\stmt{\Gamma}{m}{\kw{log}\ E(a_1 = e_1, \ldots, a_n = e_n)}{\Gamma}}
\end{mathpar}

%==============================================================================
\section{Mutability Constraints}
%==============================================================================

\fbox{$m_1 \leq m_2$} \quad (mutability ordering)

\[
\kw{pure} < \kw{view} < \kw{nonpayable} < \kw{payable}
\]

\subsection{State Access Rules}

\begin{mathpar}
\inferrule*[right=M-Pure]
  {m = \kw{pure}}
  {\text{No state variable access}}

\inferrule*[right=M-View]
  {m = \kw{view}}
  {\text{Read-only state access}}

\inferrule*[right=M-Nonpayable]
  {m \geq \kw{nonpayable}}
  {\text{Read/write state access}}

\inferrule*[right=M-Payable]
  {m = \kw{payable}}
  {\kw{msg.value} \text{ accessible}}
\end{mathpar}

\subsection{Call Mutability}

\begin{mathpar}
\inferrule*[right=M-Call]
  {\Gamma \vdash_m \mathit{caller} \\
   f \text{ has mutability } m_f \\
   m_f \leq m \lor m \geq \kw{nonpayable}}
  {\text{Call to } f \text{ is valid}}
\end{mathpar}

%==============================================================================
\section{Location Constraints}
%==============================================================================

\subsection{Type-Location Restrictions}

\begin{mathpar}
\inferrule*[right=L-HashMap]
  { }
  {\kw{HashMap} \text{ only valid in } \ell \in \{\kw{storage}, \kw{transient}\}}

\inferrule*[right=L-ExtParam]
  {f \text{ is external}}
  {\text{Parameters: } \ell = \kw{calldata}}

\inferrule*[right=L-IntParam]
  {f \text{ is internal}}
  {\text{Parameters: } \ell = \kw{memory}}

\inferrule*[right=L-Immutable]
  {\ell = \kw{code}}
  {\text{Read anywhere; write only in constructor}}
\end{mathpar}

%==============================================================================
\section{Type Constraints}
%==============================================================================

\subsection{Valid HashMap Key Types}

\[
\kappa \in \{\kw{bool}, \kw{int}N, \kw{uint}N, \kw{address}, \kw{bytes}M, \kw{Bytes}[n], \kw{String}[n], \kw{flag}\}
\]

\subsection{Valid Array Element Types}

Static arrays: $\tau$ where $\tau.\mathit{is\_valid\_element\_type} = \kw{true}$

Dynamic arrays: $\tau$ where $\tau.\mathit{\_as\_darray} = \kw{true}$

\subsection{Conversion Rules}

Valid conversions via \kw{convert}($\tau_1$, $\tau_2$):

\begin{itemize}
  \item Between integer types (with bounds checking)
  \item Between \kw{bytes}$M$ and integers (with sign extension for signed)
  \item Between \kw{address} and \kw{uint160}
  \item \kw{decimal} to/from integers (truncation toward zero)
  \item \kw{bool} from any nonzero value
  \item \kw{flag} to/from \kw{uint256} only
\end{itemize}

%==============================================================================
\section{Summary}
%==============================================================================

The Vyper type system is characterized by:

\begin{enumerate}
  \item \textbf{No implicit coercions}: Types must match exactly or through explicit subtyping
  \item \textbf{Nominal typing}: User-defined types (structs, flags, events) use identity comparison
  \item \textbf{Structural subtyping}: Bytestrings and dynamic arrays are subtypes by length
  \item \textbf{Mutability tracking}: Functions declare state access level (pure/view/nonpayable/payable)
  \item \textbf{Location tracking}: Variables have associated storage locations
  \item \textbf{Explicit conversions}: All type conversions require the \kw{convert()} builtin
  \item \textbf{Bounds checking}: Integer operations are checked for overflow at compile-time where possible
\end{enumerate}

\end{document}
